%========================================
% Section：[4] Buddhābhithuti
%========================================
% title-en: Praise of the Buddha
\section*{m4 | Buddhābhithuti}\label{sec:m04}

\begin{chantsubsection}
  \chantnote
    {\{Handa mayaṃ buddhābhithutiṃ karoma se\}}
  \chantnotetrans
    {\{Now let us recite in praise of the Buddha.\}}
  \chantgap
  \chantline
    {Yo so tathāgato arahaṃ sammāsambuddho}
    {He is the Tathāgata (the Thus-Gone One), an arahant, perfectly awakened,}

  \chantline
    {Vijjācaraṇasampanno sugato lokavidū}
    {endowed with true knowledge and conduct, the Well-gone One, knower of the worlds,}

  \chantline
    {Anuttaro purisadammasārathi}
    {unsurpassed trainer of persons to be tamed,}

  \chantline
    {Satthā devamanussānaṃ buddho bhagavā}
    {teacher of gods and humans; the awakened, Blessed One.}

  \chantline
    {Yo imaṃ lokaṃ sadevakaṃ samārakaṃ sabrahmakaṃ}
    {He has realized this world with its gods, Māras, and Brahmās,}

  \chantline
    {Sassamaṇabrāhmaṇiṃ pajaṃ sadevamanussaṃ}
    {and this generation with its ascetics and brahmins, its gods and humans,}

  \chantline
    {Sayaṃ abhiññā sacchikatvā pavedesi}
    {and teaches it after having realized it by his own direct knowledge.}

  \chantline
    {Yo dhammaṃ desesi ādikalyāṇaṃ}
    {He teaches the Dhamma that is admirable in the beginning,}

  \chantline
    {Majjhekalyāṇaṃ pariyosānakalyāṇaṃ}
    {admirable in the middle, and admirable in the end,}

  \chantline
    {Sātthaṃ sabyañjaṇaṃ kevalaparipuṇṇaṃ parisuddhaṃ}
    {with the meaning and wording complete, perfect in every respect, wholly pure.}

  \chantline
    {Brahmacariyaṃ pakāsesi}
    {He has made known the perfectly complete holy life.}

  \chantline
    {Tamahaṃ bhagavantaṃ abhipūjayāmi}
    {To that Blessed One I pay deep homage.}

  \chantline
    {Tamahaṃ bhagavantaṃ sirasā namāmi \;\;\; --- bow ---}
    {To that Blessed One I bow my head.}

  % \chantnote{--- bow ---}
\end{chantsubsection}
